% =============================================================================
%  Abstract
% =============================================================================
\chapter*{Abstract}
\addcontentsline{toc}{chapter}{Abstract}
\thispagestyle{plain}

The thesis focuses on addressing the critical issue of railway ballast
degradation, which significantly impacts the safety, stability, and
operational efficiency of railway tracks. Tailored specifically to the
unique conditions of Indian Railways, this research aims to develop a
deeper understanding of ballast behaviour under impact loading and to
predict the quantum of ballast degradation using 3D image processing as well
as machine-learning models.

At the core of the study is the modification of the European-standard
drop-weight impact test apparatus to better align with Indian railway
ballast specifications and operational conditions. The modified setup
accurately replicates on-field stress conditions experienced by ballast in
the Indian context, providing a more relevant and reliable framework for
analysis. The modified apparatus serves as a critical tool for evaluating
ballast degradation under controlled yet realistic conditions, delivering
nominal contact stresses of \SI{275}{\kilo\pascal} (at a \SI{300}{\milli\metre}
drop height) and \SI{523}{\kilo\pascal} (at a \SI{450}{\milli\metre} drop
height) corresponding to the \num{22.5}-tonne axle load characteristic of
Indian Railways.

A significant innovation in this study is the use of advanced 3D image
processing techniques to evaluate morphological changes in ballast particles
before and after impact testing. High-resolution 3D meshes are captured to
analyse changes in particle shape, size, and texture, providing a detailed
and quantitative understanding of ballast degradation. A paired-particle
protocol applied to 45 marked tracer particles across nine experimental
configurations enables precise measurement of deterioration, allowing
identification of critical thresholds and damage modes (diffuse abrasion vs.\
catastrophic fracture) that are not detectable through traditional sieve
analysis.

The outcomes of this study contribute to railroad ballast degradation
assessment using image-based morphology and machine learning, delivering a
scalable and accurate alternative to conventional physical testing, with
substantial implications for railway safety and maintenance planning.

\vspace{0.4cm}
\noindent\textbf{Keywords:} Railway ballast; drop-weight impact loading; RDSO
gradation; Fouling Index; 3D structured-light scanning; particle morphology;
under-ballast mats; ballast degradation.
